% ===============================================
% MATH 3053: Abstract Algebra I         Fall 2017
% hw_revised_cn.tex
% 中文适配版作业修订提交模板
% ===============================================
%         请仔细阅读以下说明！！！
% ===============================================
% 生成PDF提交时，将文件名改为 hwX-姓名.pdf
% X替换为作业编号，姓名替换为你的姓氏
% ===============================================


% -----------------------------------------------
% 以下序言部分可忽略，直接跳至 "START HERE" 区域
% -----------------------------------------------

\documentclass{article}
% 页面设置与基础宏包
\usepackage[margin=1in]{geometry}  % 页边距1英寸
\usepackage{amsmath,amsthm,amssymb} % 数学公式、定理环境、符号
\usepackage{ctex}   % 中文支持核心宏包（Overleaf预装）
\usepackage{pifont}                 
\renewcommand{\qedsymbol}{}
% 常用数学集合符号定义
\newcommand{\R}{\mathbb{R}}  % 实数集
\newcommand{\Z}{\mathbb{Z}}  % 整数集
\newcommand{\N}{\mathbb{N}}  % 自然数集
\newcommand{\Q}{\mathbb{Q}}  % 有理数集
\newcommand{\C}{\mathbb{C}}  % 复数集

% 定理、习题等环境定义（支持中英文标注）
\newenvironment{theorem}[2][Theorem]{\begin{trivlist}
\item[\hskip \labelsep {\bfseries #1}\hskip \labelsep {\bfseries #2.}]}{\end{trivlist}}
\newenvironment{lemma}[2][Lemma]{\begin{trivlist}
\item[\hskip \labelsep {\bfseries #1}\hskip \labelsep {\bfseries #2.}]}{\end{trivlist}}
\newenvironment{exercise}[2][Exercise]{\begin{trivlist}
\item[\hskip \labelsep {\bfseries #1}\hskip \labelsep {\bfseries #2.}]}{\end{trivlist}}
% 中文"问题"环境（与原problem并存，可按需选用）
\newenvironment{problem}[2][Problem]{\begin{trivlist}
\item[\hskip \labelsep {\bfseries #1}\hskip \labelsep {\bfseries #2.}]}{\end{trivlist}}
\newenvironment{problem_cn}[2][问题]{\begin{trivlist}
\item[\hskip \labelsep {\bfseries #1}\hskip \labelsep {\bfseries #2.}]}{\end{trivlist}}
\newenvironment{question}[2][Question]{\begin{trivlist}
\item[\hskip \labelsep {\bfseries #1}\hskip \labelsep {\bfseries #2.}]}{\end{trivlist}}
\newenvironment{corollary}[2][Corollary]{\begin{trivlist}
\item[\hskip \labelsep {\bfseries #1}\hskip \labelsep {\bfseries #2.}]}{\end{trivlist}}

% 解答环境（支持中英文标注）
\newenvironment{solution}{\begin{proof}[Solution]}{\end{proof}}
\newenvironment{solution_cn}{%
  \begin{proof}[解答]%
  \kaishu  % 切换中文为楷体（ctex 提供的命令）
}{%
  \end{proof}%
}

\begin{document}

% ------------------------------------------ %
%                 从这里开始编辑                 %
% ------------------------------------------ %

% 标题与作者信息（中文适配）
\title{《优化方法基础》课程作业\\凸集和凸函数}  % X替换为作业编号
\author{姓名：林圣翔 \quad 班级：计试2201 \quad 学号：2223312202}  % 替换括号内容

\maketitle

\begin{problem_cn}{1}
证明：范数$\boldsymbol{\|\cdot\|}$的对偶范数满足满足范数的定义：  
$$\|z\|_* = \sup \left\{ z^T x : \|x\| \leq 1 \right\} = \sup \left \{ z^T x : \|x\| = 1 \right\}$$ 
\end{problem_cn}
\begin{solution_cn} 
\quad \\1.证明： $\sup \left\{ z^T x : \|x\| \leq 1 \right\} = \sup \left \{ z^T x : \|x\| = 1 \right\}$ \\
记$A = \sup\left\{ z^T x : \|x\| \leq 1 \right\},\quad B = \sup\left\{ z^T x : \|x\| = 1 \right\}$\\
$\because \{x : \|x\| = 1 \} \subset \{x : \|x\| \leq 1\} \quad \therefore B \leq A$\\
$\forall x \in \{x : \|x\| \leq 1\}$，分两种情况：\\ \ding{172} 若$x = 0$：$z^T x = 0 \leq B$；\\
\ding{173} 若$x \neq 0$：令$x' = \frac{x}{\|x\|}$，由原范数正齐次性得$\|x'\| = 1$，则
  $ z^T x = \|x\| \cdot z^T x' \leq \|x\| \cdot B$ \\
  又$\|x\| \leq 1$，故$z^T x \leq B$。
综上，对所有$\|x\| \leq 1$均有$z^T x \leq B$，即$A \leq B$\\
2.验证范数的四个条件\\
2.1.非负性，$\forall z \in \mathbb{R}^n$，$\|z\|_* \geq 0$
\\由定义，$| z |_* = \sup { z^T x : | x | \leq 1 }$，由于 $z^T x$ 是实数，上确界存在且非负。\\ 若 $z = 0$，则 $z^T x = 0$ 对所有 $x$ 成立，故 $| z |_* = 0$。 \\若 $| z |_* = 0$，则 $z^T x = 0$ 对所有 $| x | \leq 1$ 成立。特别地，取 $x = z / | z |$（若 $z \neq 0$），则 $z^T x = | z | > 0$，矛盾。故 $z = 0$。\\
$\therefore$ $\forall z \in \mathbb{R}^n$，$\|z\|_* \geq 0$,即$\|z\|_*$满足非负性\\
2.2.正定性：若$\|z\|_* = 0$，则$z = 0$\\
必要性：若$z = 0$，则对任意$x$，$z^T x = 0$，故$\|z\|_* = \sup\{0\} = 0$。\\
充分性：若$\|z\|_* = 0$，则对所有$\|x\| \leq 1$，$z^T x \leq 0$。  \\
  假设$z \neq 0$，取$x_0 = \frac{z}{\|z\|}$（由原范数齐次性，$\|x_0\| = 1$），则$
  z^T x_0 = \frac{z^T z}{\|z\|} = \frac{\|z\|^2}{\|z\|} = \|z\| > 0 $
  与“$z^T x \leq 0$对所有$\|x\| \leq 1$成立”矛盾，故$z = 0$。\\
2.3.齐次性：$\forall z \in \mathbb{R}^n$，$\forall t \in \mathbb{R}$，$\|tz\|_* = |t| \|z\|_*$
\begin{flushleft}
$\begin{aligned}
\|tz\|_* &= \sup \left\{ (tz)^T x : \|x\| \leq 1 \right\} \\
&= \sup \left\{ t \cdot z^T x : \|x\| \leq 1 \right\} \\
&= |t| \cdot \sup \left\{ z^T x : \|x\| \leq 1 \right\} \\
&= |t| \|z\|_*
\end{aligned}$
\end{flushleft}
2.4.三角不等式：$\forall z_1, z_2 \in \mathbb{R}^n$，$\|z_1 + z_2\|_* \leq \|z_1\|_* + \|z_2\|_*$
\begin{flushleft}
$\begin{aligned}
\|z_1 + z_2\|_* &= \sup \left\{ (z_1 + z_2)^T x : \|x\| \leq 1 \right\} \\
&= \sup \left\{ z_1^T x + z_2^T x : \|x\| \leq 1 \right\}
\end{aligned}$
\end{flushleft}
对任意$\|x\| \leq 1$，由对偶范数定义：
$z_1^T x \leq \|z_1\|_*, \quad z_2^T x \leq \|z_2\|_*$\\
$\therefore$ $z_1^T x + z_2^T x \leq \|z_1\|_* + \|z_2\|_*$，即$\|z_1 + z_2\|_* \leq \|z_1\|_* + \|z_2\|_*$。\\
综上，范数$\boldsymbol{\|\cdot\|}$的对偶范数满足满足范数的定义： $\|z\|_* = \sup \left\{ z^T x : \|x\| \leq 1 \right\} = \sup \left \{ z^T x : \|x\| = 1 \right\}$ 
\end{solution_cn}
\vspace{0.25in}
\begin{problem_cn}{2}
证明： $$z^T x \leq \|x\| \|z\|_*$$  
\end{problem_cn}
\begin{solution_cn}
\quad \\$\because \quad \|z\|_* = \sup\!\{ z^T x : \|x\| \leq 1 \} = \sup\!\{ z^T x : \|x\| = 1 \} = \sup\!\left\{ z^T \frac{x}{\|x\|} : \|x\| = 1 \right\}$\\
$\therefore\quad \|z\|_* \geq z^T \frac{x}{\|x\|}$\\
$\therefore\quad z^T x \leq \|x\| \|z\|_*$ 
\end{solution_cn}
\vspace{0.25in}
\begin{problem_cn}{3}
 \( f: \mathbb{R}^n \to \mathbb{R} \),  
$$ f(x) = \log \sum_{i=1}^m \exp\left( a_i^\top x + b_i \right), $$  
其中 \( a_1, a_2, \dots, a_m \in \mathbb{R}^n \)，\( b_1, b_2, \dots, b_m \in \mathbb{R} \)。 
\begin{enumerate}
\item 求函数 \( f \) 的梯度 \( \nabla f(x) \)；  
\item 求函数 \( f \) 的二阶导数（Hessian 矩阵）\( \nabla^2 f(x) \)。  
\end{enumerate}
\end{problem_cn}
\begin{solution_cn}
\quad \\
1.设 $S(x) = \sum_{i=1}^m \exp\left(a_i^\top x + b_i\right)$,则函数可表示为$f(x) = \log S(x)$ \\
对 \( f(x) \) 求梯度，利用链式法则，有$\nabla f(x) = \frac{1}{S(x)} \nabla S(x)$,其中 $\nabla S(x) = \sum_{i=1}^m \exp\left(a_i^\top x + b_i\right) a_i $ \\
$\therefore$函数\( f \)的梯度为 $\nabla f(x) = \frac{\sum_{i=1}^m \exp\left(a_i^\top x + b_i\right) a_i}{\sum_{i=1}^m \exp\left(a_i^\top x + b_i\right)}$ \\
2.Hessian矩阵是梯度的导数：$\nabla^2 f(x) = \nabla \left( \nabla f(x) \right)$\\
由梯度表达式$\nabla f(x) = \frac{\nabla S(x)}{S(x)}$，可得：$\nabla^2 f(x) = \frac{S(x) \cdot \nabla^2 S(x) - \nabla S(x) \cdot \left( \nabla S(x) \right)^\top}{S(x)^2}$\\
$\nabla^2 S(x) = \sum_{i=1}^m \exp\left( a_i^\top x + b_i \right) \cdot a_i a_i^\top$ \\
代入得 Hessian 矩阵，可得：$ \nabla^2 f(x) = \frac{1}{S(x)^2} \left[ S(x) \cdot \sum_{i=1}^m \exp\left( a_i^\top x + b_i \right) a_i a_i^\top - \nabla S(x) \cdot \left( \nabla S(x) \right)^\top \right]$\\
利用权重函数 \( w_i(x) \)，有$\frac{\exp\left(a_i^\top x + b_i\right)}{S(x)} = w_i(x)$\\
$\therefore$函数 \( f \) 的二阶导数（Hessian 矩阵）$\nabla^2 f(x) = \sum_{i=1}^m w_i(x) \, a_i a_i^\top - \left( \sum_{i=1}^m w_i(x) \, a_i \right) \left( \sum_{i=1}^m w_i(x) \, a_i \right)^\top$

\end{solution_cn}
\vspace{0.25in}
\begin{problem_cn}{4}
令 \( A \in \mathbb{R}^{m \times n} \)，\( B \in \mathbb{R}^{p \times n} \)，\( b \in \mathbb{R}^m \)，\( d \in \mathbb{R}^p \)，集合：  
$$ P = \left\{ x \in \mathbb{R}^n : Ax \leq b \text{ 且 } Bx = d \right\} $$  
是凸集吗？为什么？
\end{problem_cn}
\begin{solution_cn}
\quad \\
1.线性不等式约束\\
由于 \( x, y \in P \)，有 \( Ax \leq b \) 和 \( Ay \leq b \)。  
则  
\begin{flushleft}
$\begin{aligned}
Az &= A\bigl(tx + (1-t)y\bigr) \\
&= t\,Ax + (1-t)\,Ay \\
&\leq t\,b + (1-t)\,b \\
&= b,
\end{aligned}
$\end{flushleft}
$\therefore$ \( z \) 满足不等式约束。\\  
2. 线性等式约束\\
由于 \( x, y \in P \)，有 \( Bx = d \) 和 \( By = d \)。  
则  
\begin{flushleft}
$\begin{aligned}
Bz &= B\bigl(tx + (1-t)y\bigr) \\
&= t\,Bx + (1-t)\,By \\
&= t\,d + (1-t)\,d \\
&= d,
\end{aligned}
$\end{flushleft}
$\therefore$ \( z \) 满足等式约束。 \\
综上，P是凸集。
\end{solution_cn}
\vspace{0.25in}
\begin{problem_cn}{5}
证明：最大值函数 \( f(x) = \max\left\{ x_1, x_2, \dots, x_n \right\} \)，\( x = [x_1, x_2, \dots, x_n]^\top \in \mathbb{R}^n \)为凸函数。
\end{problem_cn}
\begin{solution_cn}
\quad \\设 \( z = tx + (1-t)y \)，则对每个分量 \( i = 1, \dots, n \)，有$z_i = t x_i + (1-t) y_i$\\
令 \( f(x) = \max\{x_1, \dots, x_n\} = x_k \)（其中 \( x_k \) 是 \( x \) 的最大分量）， \( f(y) = \max\{y_1, \dots, y_n\} = y_l \)（其中 \( y_l \) 是 \( y \) 的最大分量）。  
此时，不等式右边为：$t f(x) + (1-t) f(y) = t x_k + (1-t) y_l$\\
考虑左边 \( f(z) \)：$f(z) = \max\{z_1, \dots, z_n\} = \max\bigl\{ t x_i + (1-t) y_i \mid i = 1, \dots, n \bigr\}$\\
$\because$对任意 \( i \)， \( x_i \leq x_k \) 且 \( y_i \leq y_l \)，$\therefore$ $t x_i + (1-t) y_i \leq t x_k + (1-t) y_l$\\
由于最大值是各分量的上确界，有：$f(z) = \max\bigl\{ t x_i + (1-t) y_i \mid i = 1, \dots, n \bigr\} \leq t x_k + (1-t) y_l$\\
代入右边的表达式，得$f(z) \leq t f(x) + (1-t) f(y)$\\
故凸函数的定义式成立\\
$\therefore$ \( f \) 是凸函数。 
\end{solution_cn}
\vspace{0.25in}
\begin{problem_cn}{6}
对于凸优化问题：  
\[
\begin{split}
\min \quad & f(x) \\
\text{s.t.} \quad & x \in X
\end{split}
\]  
如果目标函数f是可微的，可行解 \( x^* \in X \) 是最优解当且仅当 $\forall$ \( y \in X \)，有  
\[
\nabla f(x^*)^\top (y - x^*) \geq 0
\]  
\end{problem_cn}
\begin{solution_cn}
\quad \\必要性\\
假设 \( x^* \) 是最优解。对任意 \(y \in X \)，考虑函数 $\varphi(t) = f\bigl( x^* + t(y - x^*) \bigr), \quad t \in [0,1]$\\
由于 \( X \) 是凸集，对任意 \( t \in [0,1] \)，有 \( x^* + t(y - x^*) \in X \)。\\因 \( f \) 在 \( x^* \) 处可微，且 \( x^* \) 是最小值点，故 \( \varphi(t) \) 在 \( t=0 \) 处取得最小值，从而$ \varphi'(0) \geq 0$\\
$\because \varphi'(0) = \nabla f(x^*)^\top (y - x^*)$\\
$\therefore \nabla f(x^*)^\top (y - x^*) \geq 0$,必要性得证。\\
充分性\\
假设对任意 \( y \in X \)，有$\nabla f(x^*)^\top (y - x^*) \geq 0$\\
由于 \( f \) 是凸函数，对任意 \( y \in X \)，由可微凸函数的一阶性质，有$f(y) \geq f(x^*) + \nabla f(x^*)^\top (y - x^*)$\\
结合假设的不等式 \( \nabla f(x^*)^\top (y - x^*) \geq 0 \)，可得：  
$f(y) \geq f(x^*) + \nabla f(x^*)^\top (y - x^*) \geq f(x^*)$,  
即 \( f(y) \geq f(x^*) \)。
\\故 \( x^* \) 是最优解，充分性得证。\\
综上，$\forall y \in X,\ \nabla f(x^*)^\top (y - x^*) \geq 0$,结论成立

\end{solution_cn}
\vspace{0.25in}
\begin{problem_cn}{7}
证明 \( x^* = \left( 1, \dfrac{1}{2}, -1 \right)^\top \) 是如下凸优化问题的最优解：  
\[
\begin{split}
\min \quad & \dfrac{1}{2} x^\top P x + q^\top x + r \\
\text{s.t.} \quad & -1 \leq x_i \leq 1, \ i=1,2,3
\end{split}
\]  
其中：  
\[
P = \begin{bmatrix} 13 & 12 & -2 \\ 12 & 17 & 6 \\ -2 & 6 & 12 \end{bmatrix}, \quad q = \begin{bmatrix} -22 \\ -14.5 \\ 13 \end{bmatrix}, \quad r = 1
\]  
\end{problem_cn}
\begin{solution_cn}
\quad \\
二次函数 \( f(x) = \dfrac{1}{2} x^\top P x + q^\top x + r \) 的凸性等价于矩阵 \( P \) 的正定性。\\
一阶主子式：\( \Delta_1 = 13 > 0 \)\\
二阶主子式：$\Delta_2 = \begin{vmatrix} 13 & 12 \\ 12 & 17 \end{vmatrix} = 77 > 0$\\
三阶主子式 ：$\det(P) = 13 \times \begin{vmatrix} 17 & 6 \\ 6 & 12 \end{vmatrix} - 12 \times \begin{vmatrix} 12 & 6 \\ -2 & 12 \end{vmatrix} + (-2) \times \begin{vmatrix} 12 & 17 \\ -2 & 6 \end{vmatrix} = 100 > 0$\\
$\because$所有顺序主子式均为正\\$\therefore$ \( P \) 正定，\( f(x) \) 是\textbf{严格凸函数}。\\$\therefore$约束集 \( \{-1 \leq x_i \leq 1, \, i=1,2,3\} \) 为凸集，因此原问题是\textbf{严格凸优化问题}，最优解若存在则唯一。\\
$\because$ $x_1 = 1 \in [-1,1], \quad x_2 = \dfrac{1}{2} \in [-1,1], \quad x_3 = -1 \in [-1,1]$\\
$\therefore$ \( x^* \) 是可行解。  \\
$\because$
$
x^* = \begin{bmatrix} 1 \\ \dfrac{1}{2} \\ -1 \end{bmatrix}, \quad 
P = \begin{bmatrix} 13 & 12 & -2 \\ 12 & 17 & 6 \\ -2 & 6 & 12 \end{bmatrix}, \quad 
q = \begin{bmatrix} -22 \\ -14.5 \\ 13 \end{bmatrix}
$ \\
$\therefore$
$
\nabla f(x^*) = P x^* + q = = \begin{bmatrix} -1 \\ 0 \\ 2 \end{bmatrix}$\\
对于 \( x_1 = 1 \)（上界）,$\dfrac{\partial f}{\partial x_1} = -1 \leq 0$,满足上界梯度分量小于等于0的条件\\
对于 \( x_2 = \dfrac{1}{2} \)（内部），$\dfrac{\partial f}{\partial x_2} = 0$,满足内部梯度分量等于0的条件\\
对于 \( x_3 = -1 \)（下界），$\dfrac{\partial f}{\partial x_3} = 2 \geq 0$,满足下界梯度分量大于等于0的条件\\
综上，\( x^* = \left( 1, \dfrac{1}{2}, -1 \right)^\top \) 是该凸优化问题的最优解


\end{solution_cn}
% -----------------------------------------------
% 以下内容无需编辑
% -----------------------------------------------

\end{document}